\section{Redes neurais do zero}
\label{sec:nn}

\subsection{O neurônio artificial}
\begin{definicao}[Neurônio]
Um neurônio recebe números de entrada $x_1,\dots,x_n$, multiplica cada um por um
\textbf{peso} $w_i$, soma tudo com um \textbf{viés} $b$, e passa o resultado por
uma \textbf{função de ativação} $f$:
\[
y = f\!\left(\textstyle\sum_i w_i x_i + b\right).
\]
\end{definicao}

\begin{intuicao}
Os pesos dizem ``quanto cada entrada importa''. Treinar a rede é, no fundo, achar
os pesos que fazem a saída ser útil. A ativação $f$ introduz
\emph{não-linearidade}: sem ela, empilhar neurônios daria sempre uma simples
combinação linear, incapaz de representar padrões complexos. Uma ativação muito
usada é a \textbf{ReLU}, $f(z)=\max(0,z)$: deixa passar o positivo, zera o
negativo.
\end{intuicao}

\subsection{Camadas e a rede densa (MLP)}
Empilhando neurônios em \textbf{camadas} --- a saída de uma alimenta a próxima ---
obtemos uma rede. Quando cada neurônio se liga a todos da camada seguinte, é uma
camada \textbf{densa} (\emph{fully connected}); a rede toda é um \emph{Multilayer
Perceptron} (MLP).

\begin{lstlisting}[caption={Uma MLP minúscula em Keras (só para fixar a ideia).}]
from tensorflow import keras
modelo = keras.Sequential([
    keras.layers.Dense(8, activation="relu", input_shape=(3,)),
    keras.layers.Dense(1)               # saida
])
\end{lstlisting}

\subsection{Função de perda: medir o erro}
\begin{definicao}[Função de perda]
Número que mede o quão erradas estão as previsões. Treinar = ajustar os pesos
para \emph{minimizar} a perda. Para prever números, duas perdas comuns são:
\begin{align*}
\text{MSE} &= \frac{1}{n}\sum_i (\hat{y}_i - y_i)^2 \quad\text{(erro quadrático médio)},\\
\text{MAE} &= \frac{1}{n}\sum_i |\hat{y}_i - y_i| \quad\text{(erro absoluto médio)}.
\end{align*}
\end{definicao}

\begin{intuicao}[MSE vs.\ MAE]
O MSE eleva o erro ao quadrado, então pune \emph{desproporcionalmente} os
desvios grandes --- um único ponto muito fora domina a conta. O MAE trata todos
os desvios em pé de igualdade, sendo mais \textbf{robusto a \emph{outliers}}.
\end{intuicao}

\subsection{Como a rede aprende: gradiente descendente}
A perda depende dos pesos. Imagine a perda como a altitude de um terreno
montanhoso e os pesos como sua posição: queremos descer até o vale. O
\textbf{gradiente} aponta a direção de subida mais íngreme; então damos passos
na direção \emph{oposta}. O tamanho do passo é a \textbf{taxa de aprendizado}
(\emph{learning rate}).

\begin{itemize}[noitemsep]
  \item \textbf{Época} (\emph{epoch}): uma passada por todos os dados de treino.
  \item \textbf{Lote} (\emph{batch}): subconjunto processado por vez antes de
        atualizar os pesos.
  \item \textbf{Otimizador}: a receita de como dar os passos. Usamos o
        \texttt{Adam}, que ajusta o passo automaticamente e é o padrão moderno.
\end{itemize}

\subsection{Overfitting, validação e \emph{early stopping}}
\begin{definicao}[Overfitting]
Quando a rede \emph{decora} o treino (inclusive seu ruído) em vez de aprender o
padrão geral. Sinal: erro baixíssimo no treino, mas alto em dados novos.
\end{definicao}

Para detectar isso, separamos uma fatia do treino como \textbf{validação} e
monitoramos a perda nela durante o treino. Enquanto a validação melhora, a rede
aprende; quando ela começa a piorar (mesmo com o treino ainda melhorando), começou
o overfitting. O \textbf{\emph{early stopping}} interrompe o treino nesse ponto e
restaura os melhores pesos. O \textbf{\emph{dropout}} é outra defesa: durante o
treino, desliga aleatoriamente uma fração dos neurônios, forçando a rede a não
depender de detalhes específicos.

\begin{naprat}
Treinamos com a perda \textbf{MSE}, otimizador \textbf{Adam} (taxa $0{,}001$),
lotes de 32, \emph{dropout} de $0{,}2$ e \textbf{\emph{early stopping}} com
paciência de 10 épocas (restaurando os melhores pesos). Um detalhe vital para
série temporal: a validação usa o \emph{bloco final} do treino (não uma amostra
aleatória) e \texttt{shuffle=False}, para não embaralhar o tempo e recriar o
vazamento da Seção~\ref{sec:vazamento}. Na detecção (Seção~\ref{sec:deteccao})
preferimos o \textbf{MAE} como medida de erro, por sua robustez a \emph{outliers}.
\end{naprat}
